\section{Literatür verisiyle yazılım uygulamaları}
\label{sec:spss-b10}

\textbf{Amaç ve veri.} \texttt{b10.csv} ile anlamlılık, etki büyüklüğü
ve ileriye dönük örneklem planını ayırın \citep{cortez2008data}.
Gözlenen analizde $H_0:\mu=10$ ve $H_1:\mu\ne10$ sınanır.
Planlamada ise yeni bir çalışma için hedef fark 1 puandır;
bu fark, gözlenen ortalama farkından ayrı bir girdidir.

\textbf{Ortak veri.} Doğrulanan B10 uygulamasında üç yazılım da
\nolinkurl{companion/spss/csv/b10.csv} dosyasını kullanmıştır.
Dosya 395 satır ve 10 sütun içerir; kayıtların tamamı analiz edilir.
B06'nın seçim göstergeleri ve ağırlığı burada kullanılmaz.
Bu karşılaştırma Excel içe aktarmanın doğrulandığı anlamına gelmez.
Kodları \texttt{companion} klasörünü içeren paket kökünden çalıştırın;
veri sözlüğü ve hazırlık için bkz. sayfa~\pageref{sec:spss-guide}.

\subsection{IBM SPSS uygulaması}
\textbf{Başlamadan önce.} Aşağıdaki kod CSV verisini yeniden açar.
Menü yolunu kullanacaksanız veri açma ve tanımlama komutlarını
\texttt{EXECUTE} satırı dahil çalıştırın. Filtre, ağırlık ve dosya
bölme kapalı olmalıdır. \texttt{AGGREGATE}, etkin veriyi tek özet
satırıyla değiştirir; bu özeti ham veri dosyasının üzerine kaydetmeyin.
Tekrar analiz için veri açma adımından başlayın. Güç komutu SPSS
Statistics 27 ve sonrasında bulunur; eski sürümlerde yalnız bu adımı
atlayıp R veya Python planını kullanabilirsiniz.

\begin{spssadimlar}
\spssadim{Gözlenen test ve $d$ için aşağıdaki kodu kullanın. \emph{One-Sample Test} tablosunda \emph{Two-Sided p} sütununu, \texttt{LIST} çıktısında \texttt{etki\_d} değerini okuyun. Örneklem planını bu gözlenen farka göre kurmayın.}
\spssadim{\emph{Analyze $\to$ Power Analysis $\to$ Means $\to$ One-Sample T Test} yolunu açın. \emph{Estimate} için \emph{Sample size} seçin.}
\spssadim{Hedef gücü \texttt{0.80} girin. \emph{Hypothesized Values} ile \emph{Population mean}=11 ve \emph{Null value}=10 belirleyin.}
\spssadim{\emph{Population standard deviation} alanına \texttt{4.581442611} girin; \emph{Nondirectional (two-sided)} ve anlamlılık düzeyi \texttt{0.05} seçin.}
\spssadim{\emph{OK} ile gereken örneklem hacmini alın. Paylaşılan çıktıda \emph{N}=167 ve \emph{Actual Power}=0,801 görülür. Bu güç, planlanan 1 puan fark içindir; eldeki 395 kaydın gözlenen gücü değildir.}
\end{spssadimlar}

{\raggedright
\noindent\textbf{Komutların görevi.} \texttt{T-TEST} ve \texttt{COMPUTE etki\_d} gözlenen veriyi inceler. \texttt{POWER MEANS ONESAMPLE} ayrı bir ileriye dönük planlamadır; hedef güç ve alternatif ortalama önceden belirlenir.\par}

\textbf{Uygulama kodu:} \texttt{b10}. Doğrulanan veri dosyası
\texttt{companion/spss} altındaki \texttt{csv/b10.csv} dosyasıdır.

\begin{lstlisting}[style=prismSPSS]
SET DECIMAL=DOT.
GET DATA /TYPE=TXT
 /FILE='companion/spss/csv/b10.csv'
 /ENCODING='UTF8'
 /ARRANGEMENT=DELIMITED /DELCASE=LINE
 /DELIMITERS="," /QUALIFIER='"'
 /FIRSTCASE=2
 /VARIABLES=id F8.0 school F8.0 sex F8.0 age F8.0
 studytime F8.0 absences F8.0 g1 F8.0 g2 F8.0
 g3 F8.0 pass10 F8.0.
FILTER OFF.
WEIGHT OFF.
SPLIT FILE OFF.
VARIABLE LABELS g3 'Yil sonu matematik notu'.
VARIABLE LEVEL
 id school sex pass10 (NOMINAL)
 studytime (ORDINAL)
 age absences g1 g2 g3 (SCALE).
FORMATS g3 (F14.9).
EXECUTE.
\end{lstlisting}

\par\noindent\begin{minipage}{\linewidth}
\textbf{Gözlenen test ve özet.} Veri açma bloğundan sonra çalıştırın.

\begin{lstlisting}[style=prismSPSS]
T-TEST /TESTVAL=10 /VARIABLES=g3
 /CRITERIA=CI(.95).
AGGREGATE /OUTFILE=*
 /BREAK= /n=N /ort=MEAN(g3) /ss=SD(g3).
COMPUTE etki_d=(ort-10)/ss.
FORMATS n (F8.0).
FORMATS ort ss etki_d (F14.9).
LIST VARIABLES=n ort ss etki_d.
\end{lstlisting}
\end{minipage}\par

\par\noindent\begin{minipage}{\linewidth}
\textbf{Ayrı planlama adımı.} Aşağıdaki komut ham veriden yeni bir
etki tahmin etmez; belirtilen varsayımlarla gereken hacmi hesaplar.

\begin{lstlisting}[style=prismSPSS]
POWER MEANS ONESAMPLE
 /PARAMETERS TEST=NONDIRECTIONAL
 SIGNIFICANCE=.05 POWER=.80
 SD=4.581442611 MEAN=11 NULL=10.
\end{lstlisting}

\end{minipage}\par

\begin{outputguidebox}
\textbf{Paylaşılan SPSS çıktısıyla doğrulanan değerler.}
\emph{One-Sample Statistics} tablosunda $n=395$, ortalama
10,415189873, standart sapma 4,5814426110 ve standart hata
0,23051739492 görülür. Test tablosunda $t(394)=1{,}801$ ve
iki yönlü $p=0{,}072$; \texttt{LIST} çıktısında
$d=(\bar x-10)/s=0{,}090624266$ bulunmuştur.
Tek yönlü 0,036 değeri bu uygulamanın kararında kullanılmaz.

Güç tablosundaki \emph{Effect Size}=0,218 ise
$1/4{,}581442611$ planlama etkisinin yuvarlanmış değeridir;
gözlenen $d$ ile aynı değildir. \emph{Power}=0,8 hedef gücü,
\emph{Actual Power}=0,801 ise 167 kişide hesaplanan gücü gösterir.
\texttt{LIST} altındaki bir kayıt, 395 kişinin toplu özetidir;
örneklem büyüklüğünün bire düştüğü anlamına gelmez.
\end{outputguidebox}


\par\noindent\begin{minipage}{\linewidth}
\subsection{R uygulaması}
\label{sec:r-gercek-b10}

\textbf{Gözlenen analiz.} Standartlaştırılmış fark ile iki yönlü
$p$ değeri ayrı niceliklerdir. Ortalama aralığından 10 çıkarılarak
SPSS'in fark aralığı elde edilir.

\begin{lstlisting}[style=prismR]
veri <- read.csv("companion/spss/csv/b10.csv")
stopifnot(nrow(veri) == 395, ncol(veri) == 10,
  !anyNA(veri),
  all(veri$pass10 == as.integer(veri$g3 >= 10)))
sonuc <- t.test(veri$g3, mu=10,
  alternative="two.sided", conf.level=.95)
ortalama <- mean(veri$g3)
standart_sapma <- sd(veri$g3)
print(c(n=nrow(veri), ort=ortalama, ss=standart_sapma,
  sh_ort=standart_sapma/sqrt(nrow(veri)),
  t=unname(sonuc$statistic), df=unname(sonuc$parameter),
  p_iki_yonlu=sonuc$p.value, fark=ortalama-10,
  etki_d=(ortalama-10)/standart_sapma,
  ort_alt=sonuc$conf.int[1],
  ort_ust=sonuc$conf.int[2]), digits=15)
print(sum(veri$g3 == 0))
\end{lstlisting}

\textbf{Çıktıyı okuma.} Paylaşılan R çıktısında 395 satır, 10 sütun,
sıfır eksik değer ve analizde tutulan 38 sıfır not vardır.
Gözlenen etki 0,090624, iki yönlü $p$ değeri 0,072448'dir.
\end{minipage}\par

\par\noindent\begin{minipage}{\linewidth}
\textbf{İleriye dönük plan.} \texttt{strict=TRUE} iki kuyruğu da
hesaba katar; \texttt{ceiling} gereken hacmi yukarı yuvarlar.

\begin{lstlisting}[style=prismR]
plan <- power.t.test(delta=1, sd=4.581442611,
  sig.level=.05, power=.80, type="one.sample",
  alternative="two.sided", strict=TRUE)
n_plan <- ceiling(plan$n)
guc <- function(hacim) {
  power.t.test(n=hacim, delta=1, sd=4.581442611,
    sig.level=.05, type="one.sample",
    alternative="two.sided", strict=TRUE)$power
}
print(c(n_kesirli=plan$n, n_plan=n_plan,
  guc_plan=guc(n_plan), n_bir_eksik=n_plan-1,
  guc_bir_eksik=guc(n_plan-1)), digits=15)
\end{lstlisting}

\textbf{Çıktıyı okuma.} R çıktısında kesirli hacim yaklaşık
166,6755, yukarı yuvarlanan hacim 167'dir. Güç 166 kişide
0,798386, 167 kişide 0,800771 olur; hedefe aşağı yuvarlanmaz.
\end{minipage}\par

\par\noindent\begin{minipage}{\linewidth}
\subsection{Python uygulaması}
\label{sec:python-b10}

\textbf{Gözlenen analiz.} \texttt{ddof=1} örneklem standart
sapmasını seçer. \texttt{aralik} ortalamanın yüzde 95 aralığıdır.

\begin{lstlisting}[style=prismPythonApp]
import pandas as pd
import numpy as np
from scipy import stats, optimize

veri = pd.read_csv("companion/spss/csv/b10.csv")
assert veri.shape == (395, 10)
assert not veri.isna().any().any()
assert veri.pass10.eq((veri.g3 >= 10).astype(int)).all()
sonuc = stats.ttest_1samp(veri.g3, 10,
                         alternative="two-sided")
aralik = sonuc.confidence_interval(.95)
ortalama = veri.g3.mean()
standart_sapma = veri.g3.std(ddof=1)
ozet = dict(n=len(veri), ort=ortalama, ss=standart_sapma,
    sh_ort=standart_sapma/np.sqrt(len(veri)),
    t=sonuc.statistic, df=len(veri)-1,
    p_iki_yonlu=sonuc.pvalue, fark=ortalama-10,
    etki_d=(ortalama-10)/standart_sapma,
    ort_alt=aralik.low, ort_ust=aralik.high)
for ad, deger in ozet.items():
    print(f"{ad}: {deger:.13f}")
print(int((veri.g3 == 0).sum()))
\end{lstlisting}

\textbf{Çıktıyı okuma.} Paylaşılan Python çıktısı da 395 satır,
10 sütun, sıfır eksik değer ve 38 sıfır not gösterir. Gözlenen
analizin 11 özet değeri R ile 13 ondalık basamakta eşleşir.
R ve Python kodları ham veriyi tek satırlık özetle değiştirmez.
\end{minipage}\par

\par\noindent\begin{minipage}{\linewidth}
\textbf{İleriye dönük plan.} Önceki Python bloğundan sonra çalıştırın.
\texttt{nct}, merkezî olmayan $t$ dağılımıyla iki kuyruklu gücü
hesaplar. Kök çözümü yukarı yuvarlanır; bir eksik kişi de kontrol edilir.

\begin{lstlisting}[style=prismPythonApp]
def guc(hacim):
    serbestlik = hacim - 1
    kritik = stats.t.ppf(.975, serbestlik)
    kayma = np.sqrt(hacim) / 4.581442611
    return (stats.nct.cdf(-kritik, serbestlik, kayma)
            + stats.nct.sf(kritik, serbestlik, kayma))

n_kesirli = optimize.brentq(
    lambda hacim: guc(hacim)-.80, 2, 1000)
n_plan = int(np.ceil(n_kesirli))
assert guc(n_plan-1) < .80 <= guc(n_plan)
print("n_kesirli:", n_kesirli, "n_plan:", n_plan)
print("guc_plan:", guc(n_plan))
print("n_bir_eksik:", n_plan-1,
      "guc_bir_eksik:", guc(n_plan-1))
\end{lstlisting}

\textbf{Çıktıyı okuma.} Python da 167 kişi sonucunu verir.
Planlama standart sapması üç yazılımda da tam olarak
\texttt{4.581442611} girilmiştir; hesap öncesinde 4,58'e yuvarlanmaz.
\end{minipage}\par

\subsection{Sonuçların karşılaştırılması ve ortak rapor}
13 Eylül 2026'da paylaşılan SPSS tabloları ile R ve Python metin
çıktıları karşılaştırılmıştır. Gözlenen analizin 11 özet değeri
R ve Python'da 13 ondalık basamakta eşleşir; SPSS değerleri gösterilen
yuvarlama hassasiyetlerinde uyumludur. R ve Python'un kesirli hacimleri
arasında yaklaşık $1{,}1\times10^{-7}$, güçlerinde yaklaşık
$2\times10^{-12}$ düzeyinde sayısal farklar vardır. Bu farklar
167 kişilik planı değiştirmez; bütün planlama sonuçları için
13 basamakta eşitlik iddia edilmez. Proje CSV'sinden yapılan hesaplar
da sonuçları destekler. Aşağıdaki tablolar çıktıların kitap düzenine
uyarlanmış özetidir; ekran görüntüsü değildir.

\begin{table}[H]
\centering
\caption{B10'da gözlenen test ve standartlaştırılmış fark}
\label{tab:b10-dogrulanmis-gozlenen}
\begin{tabular}{@{}lr@{}}
\toprule
Ölçü & Değer \\
\midrule
Gözlem sayısı & 395 \\
Ortalama not & 10,415190 \\
Notların standart sapması & 4,581443 \\
Ortalamanın standart hatası & 0,230517 \\
Referanstan fark ($\bar x-10$) & 0,415190 \\
Gözlenen etki ($d$) & 0,090624 \\
$t$ istatistiği & 1,801122 \\
Serbestlik derecesi & 394 \\
İki yönlü $p$ değeri & 0,072448 \\
Ortalama için \%95 güven aralığı & $[9{,}961992;\ 10{,}868388]$ \\
Fark için \%95 güven aralığı & $[-0{,}038008;\ 0{,}868388]$ \\
\bottomrule
\end{tabular}
\par\smallskip
{\footnotesize Ondalıklı değerler altı basamağa yuvarlanmıştır.
Standart sapmada 394 böleni kullanılmıştır. Fark aralığı, ortalama
aralığının her iki sınırından 10 çıkarılarak elde edilir.\par}
\end{table}

\begin{table}[H]
\centering
\caption{B10'da bir puanlık hedef fark için örneklem planı}
\label{tab:b10-dogrulanmis-plan}
\begin{tabular}{@{}lr@{}}
\toprule
Planlama girdisi veya sonucu & Değer \\
\midrule
Sıfır hipotezindeki ortalama & 10 \\
Planlanan alternatif ortalama & 11 \\
Hedef fark (puan) & 1 \\
Planlama standart sapması & 4,581442611 \\
Planlama etkisi ($1/4{,}581442611$) & 0,218272 \\
İki yönlü anlamlılık düzeyi & 0,05 \\
Hedef güç & 0,80 \\
Kesirli örneklem hacmi (yaklaşık) & 166,6755 \\
Gereken en küçük tam sayı hacim & 167 \\
166 kişide hesaplanan güç & 0,798386 \\
167 kişide hesaplanan güç & 0,800771 \\
\bottomrule
\end{tabular}
\par\smallskip
{\footnotesize Kesirli hacim R ve Python hesaplarından gelir.
SPSS doğrudan 167 kişi ve üç basamakla 0,801 güç göstermiştir.
Güç hesabı bağımsız normal gözlemler modelinde merkezî olmayan
$t$ dağılımının iki kuyruğunu kullanır.\par}
\end{table}

\Needspace{13\baselineskip}
\begin{outputguidebox}
Tablo~\ref{tab:b10-dogrulanmis-gozlenen}, eldeki 395 kaydı inceler.
İki yönlü $p=0{,}072448>0{,}05$ olduğundan $H_0$ reddedilmez;
bu, ortalamanın kesinlikle 10 olduğunun veya farkın eğitim bakımından
önemsiz olduğunun kanıtı değildir. Sıfır notlu 38 kayıt korunmuştur.

Tablo~\ref{tab:b10-dogrulanmis-plan} ise farklı bir soruya yanıt
verir: 1 puanlık gerçek farkı, belirtilen standart sapma ve model
altında yüzde 80 güçle saptamak için kaç kişi gerekir?
Gözlenen $d=0{,}090624$ yerine planlama etkisi 0,218272 kullanılır.
167 kişideki 0,800771 güç, gözlenen $p$ değerinden türetilmiş
geriye dönük güç değildir; 395 kaydın 167'sini seçme önerisi de değildir.

Bir puanın anlamlı hedef sayılması ve planlama standart sapmasının
yeni çalışmaya uygunluğu ayrıca gerekçelendirilmelidir. Kayıplar ve
kümelenme bu hesaba dahil değildir. İki okuldan gelen veride yazılımların
uyumu, bağımsızlığı veya daha geniş bir evrene genellenebilirliği kanıtlamaz.
\end{outputguidebox}

\Needspace{10\baselineskip}
\textbf{Raporlama örneği (doğrulanmış çıktılarla).}
``395 kayıtta ortalama not 10,415190 ve standart sapma 4,581443
bulunmuştur. 10 puan referansına karşı iki yönlü testte
$t(394)=1{,}801122$, $p=0{,}072448$ ve standartlaştırılmış fark
$d=0{,}090624$ elde edilmiştir. Sıfır hipotezi yüzde 5 düzeyinde
reddedilmemiştir. Ayrı bir ileriye dönük planlamada, 1 puanlık hedef
fark, 4,581442611 standart sapma, iki yönlü 0,05 anlamlılık düzeyi
ve 0,80 hedef güç için en az 167 kişi gerektiği hesaplanmıştır;
bu hacimde güç 0,800771'dir. Plan bağımsız normal gözlemler modeline
dayanır; kayıp ve kümelenme etkilerini içermez. İki okulla sınırlı
kaynak verinin temsil gücü bu hesaplarla doğrulanmış değildir.''


\textbf{Görev.} Bir puanlık farkın eğitim bakımından neden önemli kabul
edilebileceğini gerekçelendirin. Bu gerekçe yoksa yalnız güç hesabı yapmak
çalışmanın bilimsel önemini sağlamaz. Gözlenen $p$ değerinden geriye dönük
güç üretip sonucu yeniden değerlendirmeyin. 166 kişiye aşağı yuvarlamanın
neden hedefi karşılamadığını açıklayın. SPSS'te sonraki uygulamaya
geçerken tek satırlık özet yerine gerekli ham veriyi yeniden açın.